\section{Теорема о взаимнооднозначном соответствии и об изоморфизме между $L(V,V)$ и пространством квадратных матриц $n$-го порядка}

\begin{theorem}
Пусть $V$ --- $n$-мерное линейное пространство над полем $F$, и в нем зафиксирован базис $E$. Отображение, сопоставляющее каждому линейному оператору $\varphi \in L(V, V)$ его матрицу $A_\varphi$ в базисе $E$, является изоморфизмом алгебры линейных операторов $L(V, V)$ и алгебры квадратных матриц $M_n(F)$.
\end{theorem}

\begin{proof}
Необходимо доказать, что отображение $\varphi \mapsto A_\varphi$ является биекцией и сохраняет алгебраические операции.

\begin{enumerate}
    \item \textbf{Биективность:} 
    По теореме о задании оператора на базисе, для любой матрицы $A \in M_n(F)$ существует единственный оператор $\varphi$, столбцы матрицы которого в базисе $E$ совпадают со столбцами $A$. Обратно, каждому оператору соответствует ровно одна матрица. Следовательно, соответствие взаимно однозначно.

    \item \textbf{Сохранение сложения и умножения на скаляр:} 
    Пусть $\varphi, \psi \in L(V, V)$ имеют матрицы $A$ и $B$. Тогда $\varphi(e_j) = \sum_i a_{ij} e_i$ и $\psi(e_j) = \sum_i b_{ij} e_i$. 
    Для суммы: $(\varphi + \psi)(e_j) = \varphi(e_j) + \psi(e_j) = \sum_i (a_{ij} + b_{ij}) e_i$, что соответствует матрице $A + B$.
    Для скаляра: $(\lambda \varphi)(e_j) = \lambda \varphi(e_j) = \sum_i (\lambda a_{ij}) e_i$, что соответствует матрице $\lambda A$.

    \item \textbf{Сохранение умножения (композиции):} 
    Пусть $\varphi$ имеет матрицу $A = (a_{ij})$, а $\psi$ --- матрицу $B = (b_{ij})$. Найдем матрицу композиции $\varphi \circ \psi$. 
    По определению матрицы оператора $\psi$:
    \[
    \psi(e_j) = \sum_{k=1}^n b_{kj} e_k.
    \]
    Применим оператор $\varphi$ к обеим частям этого равенства. В силу линейности $\varphi$:
    \[
    (\varphi \circ \psi)(e_j) = \varphi\left(\sum_{k=1}^n b_{kj} e_k\right) = \sum_{k=1}^n b_{kj} \varphi(e_k).
    \]
    Теперь подставим разложение $\varphi(e_k) = \sum_{i=1}^n a_{ik} e_i$:
    \[
    (\varphi \circ \psi)(e_j) = \sum_{k=1}^n b_{kj} \left( \sum_{i=1}^n a_{ik} e_i \right) = \sum_{i=1}^n \left( \sum_{k=1}^n a_{ik} b_{kj} \right) e_i.
    \]
    Коэффициент при $e_i$ в этом разложении равен $\sum_{k=1}^n a_{ik} b_{kj}$, что в точности является элементом $(i, j)$ матричного произведения $A \cdot B$. Следовательно, матрицей оператора $\varphi \circ \psi$ является произведение матриц $A_\varphi A_\psi$.
\end{enumerate}
Таким образом, все условия изоморфизма алгебр выполнены.
\hfill$\blacksquare$
\end{proof}

\begin{corollary}
Поскольку пространство квадратных матриц $M_n(F)$ имеет размерность $n^2$, то и пространство линейных операторов $L(V, V)$ также имеет размерность $n^2$.
\end{corollary}
